% Intended LaTeX compiler: pdflatex
\documentclass[a4paper,12pt]{article}
\usepackage[utf8]{inputenc}
\usepackage[T1]{fontenc}
\usepackage{graphicx}
\usepackage{longtable}
\usepackage{wrapfig}
\usepackage{rotating}
\usepackage[normalem]{ulem}
\usepackage{amsmath}
\usepackage{amssymb}
\usepackage{capt-of}
\usepackage{hyperref}
\usepackage{\string~/.emacs.d/common}
\author{mahmood}
\date{\today}
\title{weighted graph}
\hypersetup{
 pdfauthor={mahmood},
 pdftitle={weighted graph},
 pdfkeywords={},
 pdfsubject={},
 pdfcreator={Emacs 28.2 (Org mode 9.5.5)}, 
 pdflang={English}}
\begin{document}

\maketitle
\begin{definition}
a \textbf{weighted graph} is \(G=(V,E,w)\) such that \(G=(V,E)\) is a \href{20230314001051-graph.org}{graph} (directed or not) and \(w: E \to R\) is a weight function that gives each edge \(e \in E\) a value \(w(e)\). this value is called the \textbf{weight} or \textbf{length} of the edge \(e\)\\
\end{definition}
\begin{definition}
let \(G=(V,E,w)\) be a weighted graph, let \(\alpha=u_0,u_1,\dots,u_l\) be a \textbf{path} in \(G\), the \textbf{length} of the path \(\alpha\) is defined as the sum of weights of the edges in the path\\
$\backslash$[\\
\begin{tabular}{ll}
\(\alpha\) & =\(\sum\)\textsubscript{i=0}\textsuperscript{l-1} w(u\textsubscript{i,u}\textsubscript{i+1})\\
\end{tabular}
$\backslash$]\\
\end{definition}
\end{document}